% ch2.tex -- Module 6, Chapter 2: The Motor Effect
% Syllabus: NSW Physics Stage 6 (2017), Module 6, IQ2 (M6.IQ2.01-03); Year 11 field refresher.
% Directions checked with vectors: +x right, +y up the page, +z out of the page;
% force on a wire = I l x B; field of a wire = right-hand grip rule.
% Source error fixed here: 7c p3 draws the S-face loop with an anticlockwise current
% (it must be clockwise), see Figure ch2:fig:fields(c).
\tikzset{
  chTwocurrent/.style={->, line width=1.4pt, color=hscink},
  chTwowire/.style={line width=2.2pt, color=hscink!70},
  chTwoloop/.style={line width=1.1pt, color=hscink, midarrow},
  chTwowireout/.pic={\draw[line width=0.9pt, fill=white] (0,0) circle (0.2);
    \fill (0,0) circle (0.06);},
  chTwowirein/.pic={\draw[line width=0.9pt, fill=white] (0,0) circle (0.2);
    \draw[line width=0.9pt] (-0.13,-0.13) -- (0.13,0.13) (-0.13,0.13) -- (0.13,-0.13);},
  chTwofout/.pic={\fill[hscfield] (0,0) circle (0.055);},
  chTwofin/.pic={\draw[hscfield, line width=0.75pt] (-0.085,-0.085) -- (0.085,0.085)
    (-0.085,0.085) -- (0.085,-0.085);},
  chTwoFout/.pic={\draw[hscforce, line width=1.1pt, fill=white] (0,0) circle (0.2);
    \fill[hscforce] (0,0) circle (0.06);},
  chTwoarc/.style={line width=0.8pt, color=hscfield, postaction={decorate},
    decoration={markings, mark=at position 0.15 with {\arrow{Stealth[length=2.2mm, width=1.7mm]}}}},
  chTwoFin/.pic={\draw[hscforce, line width=1.1pt, fill=white] (0,0) circle (0.2);
    \draw[hscforce, line width=1.1pt] (-0.13,-0.13) -- (0.13,0.13) (-0.13,0.13) -- (0.13,-0.13);},
}

\chapter{The Motor Effect}\label{ch2:chap}
\covers{M6.IQ2.01, M6.IQ2.02, M6.IQ2.03}

\inquiry{Under what circumstances is a force produced on a current-carrying conductor in a
magnetic field?}

\begin{goals}
  \item describe the field of a straight wire, a loop and a solenoid, and use
        $B = \mu_0 I/(2\pi r)$ and $B = \mu_0 NI/L$
  \item use $F = lI_\perp B = lIB\sin\theta$ and the right-hand palm rule, and state when the
        force is largest and when it is zero
  \item explain why parallel wires attract or repel, and use $F/l = (\mu_0/2\pi)(I_1 I_2/r)$
  \item link the SI definition of the ampere to Newton's third law
  \item describe and evaluate the two key investigations
\end{goals}

%% ==================================================================
\section{Magnetic fields made by currents (Year 11 refresher)}\label{ch2:sec:fields}

\begin{bigidea}
Every electric current makes a magnetic field around it. Around a straight wire the field lines
are circles. A loop or a solenoid makes a field like a bar magnet's, with an N end and an S end.
\end{bigidea}

\begin{keywords}
  \kw{Magnetic field ($B$)}{a region where a magnet or a moving charge feels a force.}
  \kw{Conventional current}{the way positive charge would flow (out of the $+$ terminal,
    round the circuit, back to the $-$ terminal).}
  \kw{Right-hand grip rule}{thumb along the current; curled fingers show the field.}
  \kw{Solenoid}{a long coil of wire with many turns.}
  \kw{Uniform field}{the same strength and direction everywhere.}
\end{keywords}

\begin{refresher}{Field lines and the dot--cross code}
\term{Field lines} show the way a compass needle's N end would point. Closer lines mean a stronger
field. On a page, a dot ($\bullet$) means out of the page (an arrow tip coming at you). A cross
($\times$) means into the page (an arrow tail going away).
\end{refresher}

A current makes its own \term{magnetic field}: a region where a magnet or a moving charge feels a
force. Hans Ørsted found this in 1820, when a current made a compass needle swing. Field strength
$B$ is measured in \term{tesla} (T).

Every hand rule uses the \term{conventional current}: the way positive charge would flow. It
leaves the $+$ terminal of the battery, goes round the circuit, and comes back to the $-$
terminal. In a metal wire, electrons really move, and they flow the opposite way.

\textbf{Straight wire.} The field lines are circles around the wire (\cref{ch2:fig:fields}(a)
and (b)). The circles get further apart as you move away from the wire. This shows that the field
is weaker there.

\begin{figure}
  \centering
  \begin{tikzpicture}[x=1cm, y=1cm]
    \tikzset{chTwolab/.style={diagramlabel, font=\footnotesize, align=center}}
    % (a) straight wire, end-on
    \begin{scope}
      \foreach \r in {0.35,0.75,1.35} {\draw[field] (\r,0) arc (0:360:\r);}
      \pic at (0,0) {chTwowireout};
      \node[diagramlabel, text=hscfield] at (1.25,1.2) {$B$};
      \node[chTwolab] at (0,-1.75) {(a) wire, end-on:\\ $I$ out of the page};
    \end{scope}
    % (b) straight wire, side view
    \begin{scope}[xshift=3.2cm]
      \draw[chTwowire] (0,-1.3) -- (0,1.3);
      \draw[chTwocurrent] (0,-0.55) -- (0,0.55);
      \node[diagramlabel] at (0,1.52) {$I$};
      \foreach \y in {-1.1,-0.55,0,0.55,1.1} {\pic at (-0.32,\y) {chTwofout}; \pic at (0.32,\y) {chTwofin};}
      \foreach \y in {-0.85,-0.28,0.28,0.85} {\pic at (-0.7,\y) {chTwofout}; \pic at (0.7,\y) {chTwofin};}
      \foreach \y in {-0.6,0.6} {\pic at (-1.08,\y) {chTwofout}; \pic at (1.08,\y) {chTwofin};}
      \node[chTwolab] at (0,-1.75) {(b) wire, side view:\\ $I$ up the page};
    \end{scope}
    % (c) loops, face-on: anticlockwise = N face, clockwise = S face
    \begin{scope}[xshift=6.35cm]
      \draw[chTwoloop] (-0.15,0.1) ++(270:0.58) arc (270:630:0.58);
      \node[diagramlabel, font=\footnotesize, anchor=south] at (-0.15,0.84) {$I$};
      \foreach \p in {(-0.15,0.1),(-0.38,0.1),(0.08,0.1),(-0.15,0.33),(-0.15,-0.13)} {\pic at \p {chTwofout};}
      \node[diagramlabel, font=\footnotesize\bfseries, text=hscforce] at (-0.15,-0.78) {N face};
      \begin{scope}[xshift=1.45cm]
        \draw[chTwoloop] (-0.15,0.1) ++(270:0.58) arc (270:-90:0.58);
        \node[diagramlabel, font=\footnotesize, anchor=south] at (-0.15,0.84) {$I$};
        \foreach \p in {(-0.15,0.1),(-0.38,0.1),(0.08,0.1),(-0.15,0.33),(-0.15,-0.13)} {\pic at \p {chTwofin};}
        \node[diagramlabel, font=\footnotesize\bfseries, text=hscvel] at (-0.15,-0.78) {S face};
      \end{scope}
      \node[chTwolab] at (0.58,-1.75) {(c) loops, face-on:\\ N face and S face};
    \end{scope}
    % (d) solenoid cut along its axis
    \begin{scope}[xshift=10.95cm]
      \fill[hscbrand!6] (-1.5,-0.42) rectangle (1.5,0.42);
      \foreach \x in {-1.25,-0.75,...,1.3} {\pic at (\x,0.42) {chTwowireout}; \pic at (\x,-0.42) {chTwowirein};}
      \foreach \s in {1,-1} {
        \draw[hscfield, line width=0.8pt, postaction={decorate},
          decoration={markings, mark=at position 0.16 with {\arrow{Stealth[length=2.2mm, width=1.7mm]}},
          mark=at position 0.64 with {\arrow{Stealth[length=2.2mm, width=1.7mm]}}}]
          (-1.6,0.17*\s) -- (1.6,0.17*\s) .. controls (2.5,0.25*\s) and (2.45,1.3*\s) .. (0,1.3*\s)
          .. controls (-2.45,1.3*\s) and (-2.5,0.25*\s) .. cycle;
      }
      \node[font=\small\bfseries, text=hscforce] at (1.78,0) {N};
      \node[font=\small\bfseries, text=hscvel] at (-1.78,0) {S};
      \node[chTwolab] at (0,-1.75) {(d) solenoid, cut in half:\\ $I$ out on top, in below};
    \end{scope}
  \end{tikzpicture}
  \caption{Fields made by currents. (a), (b) A long straight wire: the field lines are circles
    round the wire. The field is weaker further away (circles further apart in (a), fewer dots and
    crosses in (b)). (c) A loop acts like a flat magnet. The arrow on each loop shows the current.
    (d) A solenoid cut in half along its length. Inside, the field is uniform and goes from S to N.
    Outside, it looks like a bar magnet's field.}\label{ch2:fig:fields}
\end{figure}

To find the direction of the field, use the \term{right-hand grip rule} (\cref{ch2:fig:grip}).
It links the way the current flows to the way the field lines go round the wire.

\begin{sidebyside}[0.40]
  \centering
  \begin{tikzpicture}[x=1cm, y=1cm]
    \tikzset{hand/.style={draw=hscink, line width=0.8pt, fill=hscbrand!7}}
    % back half of a field line (behind the wire), then the wire
    \draw[hscfield, line width=0.9pt, postaction={decorate}, decoration={markings,
      mark=at position 0.3 with {\arrow{Stealth[length=2.2mm, width=1.7mm]}}}]
      (0.95,-1.05) arc (0:180:0.95 and 0.3);
    \draw[chTwowire] (0,-1.7) -- (0,1.75);
    % front half of the field line, with its arrows (away from the wire)
    \draw[hscfield, line width=0.9pt, postaction={decorate}, decoration={markings,
      mark=at position 0.3 with {\arrow{Stealth[length=2.2mm, width=1.7mm]}},
      mark=at position 0.75 with {\arrow{Stealth[length=2.2mm, width=1.7mm]}}}]
      (-0.95,-1.05) arc (180:360:0.95 and 0.3);
    % fist (curled fingers) and thumb
    \draw[hand, rounded corners=8pt] (-0.6,-0.55) rectangle (0.6,0.5);
    \foreach \y in {-0.29,-0.03,0.23} {\draw[hscink, line width=0.6pt] (-0.6,\y) .. controls (0,\y-0.1) .. (0.6,\y);}
    \draw[hand, rounded corners=5pt] (-0.17,0.32) rectangle (0.17,1.25);
    \draw[chTwocurrent] (0,1.3) -- (0,1.95) node[above, inner sep=2pt] {$I$};
    \node[diagramlabel, anchor=east, align=right] at (-0.3,0.9) {right\\ hand};
    \node[diagramlabel, align=left, anchor=west] at (0.32,0.95) {thumb:\\ current $I$};
    \node[diagramlabel, align=left, anchor=west, text=hscfield] at (1.0,-1.05) {curled fingers:\\ field $B$};
  \end{tikzpicture}
  \captionof{figure}{The right-hand grip rule.}\label{ch2:fig:grip}
\nextside
  \begin{enumerate}
    \item Hold the wire in your right hand.
    \item Point your \textbf{thumb} along the conventional current $I$.
    \item Your curled \textbf{fingers} show the way the field lines go round the wire.
  \end{enumerate}
  Check with \cref{ch2:fig:fields}(a). The current comes out of the page, so the field lines go
  anticlockwise.
\end{sidebyside}

\begin{toolkit}{Prefixes, rearranging and the sign $\propto$}
\textbf{Prefixes:} change to SI units first.\\
c $= 10^{-2}$: $\qty{3.0}{cm} = \qty{0.030}{m}$.\\
m $= 10^{-3}$: $\qty{45}{mT} = \qty{0.045}{T}$; \ $\qty{600}{mA} = \qty{0.600}{A}$; \
$\qty{5.0}{mm} = \qty{5.0e-3}{m}$.\\
\textmu\ $= 10^{-6}$: $\qty{82}{\micro T} = \qty{8.2e-5}{T}$.\\
Grams to kilograms: divide by 1000, so $\qty{0.35}{g} = \qty{3.5e-4}{kg}$.\\
\textbf{Rearranging:} do the same to both sides. Start with $B = \mu_0 I/(2\pi r)$. Multiply both
sides by $r$: $Br = \mu_0 I/(2\pi)$. Then divide both sides by $B$: $r = \mu_0 I/(2\pi B)$.\\
\textbf{The sign $\propto$} means ``is proportional to''. $B \propto 1/r$: double $r$, and $B$ halves.
\end{toolkit}

\textbf{Current loop.} Bend a wire into a circle to make a \term{current loop}: one turn of wire.
Every part of the loop makes the field point the same way through the middle. So a loop acts like
a flat magnet with an N face and an S face (\cref{ch2:fig:fields}(c)). Curl your right fingers the
way the current goes round. Your thumb points out of the N face. Seen face-on, an anticlockwise
current gives an N face, and a clockwise current gives an S face.

\textbf{Solenoid.} A \term{solenoid} is a long coil: many loops side by side. Their fields add up.
Inside, the field is strong and \term{uniform}: the same strength and direction everywhere
(\cref{ch2:fig:fields}(d)). Outside, it looks like a bar magnet's field. To find the N end, use the
same right-hand rule as for a loop. Curl your right fingers the way the current goes round the
turns. Your thumb points to the N end.

The formulas below use $\mu_0$ (say ``mew-nought''), the \term{magnetic permeability constant}.
It is a fixed number of nature: $\mu_0 = 4\pi \times 10^{-7}$ \si{N.A^{-2}}. It sets how strong
a field a current makes in empty space. You will find it on the data sheet.

\begin{formulacard}{Field of a long straight wire}
  \cardformula{B = \frac{\mu_0 I}{2\pi r}}
  \begin{symbols}
    \sym{$B$}{magnetic field strength at distance $r$}{\si{T}}
    \sym{$\mu_0$}{magnetic permeability constant, $4\pi \times 10^{-7}$}{\si{N.A^{-2}}}
    \sym{$I$}{current in the wire}{\si{A}}
    \sym{$r$}{shortest distance from the centre of the wire}{\si{m}}
  \end{symbols}
  \inwords{a bigger current gives a stronger field. Double the distance and the field halves.}
  \whentouse{near a long, straight wire, away from its ends.}
  \whennottouse{for a coil, or to find the direction (use the grip rule).}
  \rearranged{$I = \dfrac{2\pi r B}{\mu_0}$ \qquad $r = \dfrac{\mu_0 I}{2\pi B}$ \qquad
    shortcut: $\dfrac{\mu_0}{2\pi} = \qty{2e-7}{N.A^{-2}}$}
  \tensecond{$I = \qty{5.0}{A}$, $r = \qty{0.10}{m}$:
    $B = \dfrac{2 \times 10^{-7} \times 5.0}{0.10} = \qty{1.0e-5}{T}$.}
  \trap{$B \propto \dfrac{1}{r}$, not $\dfrac{1}{r^2}$. Measure $r$ in metres, from the centre of the wire.}
\end{formulacard}

\begin{formulacard}{Field inside a solenoid}
  \cardformula{B = \frac{\mu_0 NI}{L}}
  \begin{symbols}
    \sym{$B$}{magnetic field strength inside the solenoid}{\si{T}}
    \sym{$\mu_0$}{magnetic permeability constant, $4\pi \times 10^{-7}$}{\si{N.A^{-2}}}
    \sym{$N$}{number of turns}{no unit}
    \sym{$I$}{current}{\si{A}}
    \sym{$L$}{length of the solenoid}{\si{m}}
  \end{symbols}
  \inwords{more turns, more current or a shorter coil gives a stronger field inside.}
  \whentouse{inside a long solenoid with an air core. The field there is uniform.}
  \whennottouse{outside the solenoid; for one straight wire; or for a coil with an iron core (the
    field is then much stronger than the formula gives).}
  \rearranged{$N = \dfrac{BL}{\mu_0 I}$ \qquad $I = \dfrac{BL}{\mu_0 N}$ \qquad $L = \dfrac{\mu_0 NI}{B}$}
  \tensecond{$N = 200$, $I = \qty{2.0}{A}$, $L = \qty{0.10}{m}$:
    $B = \dfrac{4\pi \times 10^{-7} \times 200 \times 2.0}{0.10} = \qty{5.0e-3}{T}$.}
  \trap{capital $L$ is the length of the coil, not the length of wire used to wind it. Put $L$ in
    metres.}
\end{formulacard}

\begin{workedexample}{Field beside a long wire}\label{ch2:we:wirefield}
A long straight wire lies in the page and carries \qty{12}{A} up the page. Point P is
\qty{3.0}{cm} to the right of the wire. Find the size and direction of the field at P.
\given{$I = \qty{12}{A}$, \ $r = \qty{3.0}{cm} = \qty{0.030}{m}$}
\find{the size and direction of $B$ at P}
\formula{$B = \dfrac{\mu_0 I}{2\pi r}$, with $\dfrac{\mu_0}{2\pi} = \qty{2e-7}{N.A^{-2}}$; then
  the right-hand grip rule}
\substitute{\eqlines{B &= \frac{2 \times 10^{-7}\,\unit{N.A^{-2}} \times \qty{12}{A}}{\qty{0.030}{m}}
  = \qty{8.0e-5}{T}}

  Direction: right thumb up the page. On the right of the wire, the curled fingers go into the
  page (\cref{ch2:fig:fields}(b)).}
\answer{\final{$B = \qty{8.0e-5}{T}$ into the page}}
\checkstep{units: $\unit{N.A^{-2}} \times \unit{A} \div \unit{m} = \unit{N.A^{-1}.m^{-1}}$, and
  $\qty{1}{T} = \qty{1}{N.A^{-1}.m^{-1}}$ (you will see why in \cref{ch2:sec:motor}) \cmark. The
  field is close to the size of Earth's field (about \qty{5e-5}{T}).}
\end{workedexample}

\begin{workedexample}{The field inside a solenoid}\label{ch2:we:solenoid}
A solenoid \qty{20}{cm} long has 400 turns and carries \qty{1.5}{A}. You look at its right-hand
end. From there, the current goes anticlockwise round the turns. Find the size of the field
inside, and say which end is N.
\given{$N = 400$, \ $I = \qty{1.5}{A}$, \ $L = \qty{20}{cm} = \qty{0.20}{m}$. Seen from the
  right-hand end, the current is anticlockwise.}
\find{the size of $B$ inside; which end is N}
\formula{$B = \dfrac{\mu_0 NI}{L}$; then the right-hand rule for a coil}
\substitute{\eqlines{B &= \frac{4\pi \times 10^{-7}\,\unit{N.A^{-2}} \times 400 \times
  \qty{1.5}{A}}{\qty{0.20}{m}} = \qty{3.77e-3}{T}}

  Direction: curl your right fingers anticlockwise, the way you see the current from the
  right-hand end. Your thumb points towards you, out of that end. So the right-hand end is N.
  Inside, the field goes from S to N, so it points to the right.}
\answer{\final{$B = \qty{3.8e-3}{T}$}, pointing to the right inside; the right-hand end is N.}
\checkstep{units: $\unit{N.A^{-2}} \times \unit{A} \div \unit{m} = \unit{N.A^{-1}.m^{-1}} =
  \unit{T}$ \cmark. Seen from an end, an anticlockwise current always means N
  (\cref{ch2:fig:fields}(c)) \cmark.}
\end{workedexample}

\begin{workedexample}{Adding the fields of two wires}\label{ch2:we:twofields}
\begin{sidebyside}[0.50]
Two long straight wires lie side by side, \qty{10.0}{cm} apart (\cref{ch2:fig:twofields}). Wire P (on the left)
carries \qty{6.0}{A} out of the page. Wire Q (on the right) carries \qty{4.0}{A} into the page.
Find the field at M, halfway between them.
\nextside
  \centering
  \begin{tikzpicture}[x=1cm, y=1cm]
    % small field lines round each wire: anticlockwise round P, clockwise round Q
    \draw[field] (0,0) ++(180:0.55) arc (180:540:0.55);
    \draw[field] (4,0) ++(0:0.55) arc (0:-360:0.55);
    \pic at (0,0) {chTwowireout}; \pic at (4,0) {chTwowirein};
    \node[diagramlabel, anchor=south] at (0,0.62) {P: \qty{6.0}{A}, out};
    \node[diagramlabel, anchor=south] at (4,0.62) {Q: \qty{4.0}{A}, in};
    % point M and the two fields there
    \fill (2,0) circle (1.6pt);
    \node[diagramlabel, anchor=north] at (2,-0.08) {M};
    \draw[->, line width=1.2pt, hscfield] (1.82,0.12) -- (1.82,1.32)
      node[left, inner sep=2pt] {$B_\mathrm{P}$};
    \draw[->, line width=1.2pt, hscfield] (2.18,0.12) -- (2.18,0.92)
      node[right, inner sep=2pt] {$B_\mathrm{Q}$};
    % distances
    \draw[<->, hscmuted, line width=0.5pt] (0,-0.85) -- (4,-0.85);
    \node[diagramlabel, fill=white, inner sep=1pt] at (2,-0.85) {\qty{10.0}{cm}};
  \end{tikzpicture}
  \captionof{figure}{End-on view. Both fields point up at M.}\label{ch2:fig:twofields}
\end{sidebyside}
\given{$I_\mathrm{P} = \qty{6.0}{A}$, \ $I_\mathrm{Q} = \qty{4.0}{A}$, \ M is
  $r = \qty{5.0}{cm} = \qty{0.050}{m}$ from each wire}
\find{the size and direction of the total field at M}
\formula{$B = \dfrac{\mu_0 I}{2\pi r}$ for each wire. Add the fields if they point the same way;
  subtract if they point opposite ways.}
\substitute{\eqlines{B_\mathrm{P} &= \frac{2 \times 10^{-7}\,\unit{N.A^{-2}} \times \qty{6.0}{A}}{\qty{0.050}{m}}
  = \qty{2.4e-5}{T}\\
  B_\mathrm{Q} &= \frac{2 \times 10^{-7}\,\unit{N.A^{-2}} \times \qty{4.0}{A}}{\qty{0.050}{m}}
  = \qty{1.6e-5}{T}}

  Directions (grip rule, see \cref{ch2:fig:twofields}): P's current is out of the page, so its
  field lines go anticlockwise round P. On the right-hand side of an anticlockwise circle, the
  line points up. M is on the right of P, so at M, $B_\mathrm{P}$ points up.

  Q's current is into the page, so its field lines go clockwise round Q. On the left-hand side of
  a clockwise circle, the line also points up. M is on the left of Q, so at M, $B_\mathrm{Q}$
  points up too.

  Both fields point up, so add them:
  $B = \qty{2.4e-5}{T} + \qty{1.6e-5}{T} = \qty{4.0e-5}{T}$.}
\answer{\final{$B = \qty{4.0e-5}{T}$ towards the top of the page}}
\checkstep{between opposite currents the two fields add. If both currents went the same way, the
  fields would subtract, giving only \qty{8.0e-6}{T}.}
\end{workedexample}

\begin{mistakes}
\begin{itemize}
  \item \xmark\ $B \propto \dfrac{1}{r^2}$. \cmark\ For a wire, $B \propto \dfrac{1}{r}$: double
        the distance and the field halves.
  \item \xmark\ Thumb along the electron flow. \cmark\ Thumb along the conventional current.
  \item \xmark\ Inside a solenoid the field goes from N to S. \cmark\ Inside it goes from S to N;
        outside it goes from N to S.
\end{itemize}
\end{mistakes}

\begin{examtip}
``Draw the field of this solenoid'' (3 marks). One mark each for:
\begin{enumerate}
  \item arrows the right way (inside from S to N, outside from N to S);
  \item closed loops that are straight and evenly spaced inside;
  \item N and S at the correct ends.
\end{enumerate}
``Find the field at P'' (two or more wires):
\begin{enumerate}
  \item find the direction of each wire's field at P (grip rule);
  \item add the fields that point the same way; subtract the fields that point opposite ways.
\end{enumerate}
\end{examtip}

\begin{quickcheck}
  \item A long straight wire carries \qty{8.0}{A}. Calculate the field \qty{4.0}{cm} from it.
  \item The field \qty{2.0}{cm} from a long wire is \qty{6.0e-5}{T}. Find the field
        \qty{6.0}{cm} from the wire.
  \item A solenoid \qty{25}{cm} long has 500 turns and carries \qty{1.2}{A}. Calculate the field
        inside it.
  \item You look at one end of a solenoid. The current goes clockwise round the turns. Is this
        end N or S?
\end{quickcheck}

%% ==================================================================
\section{The motor effect: force on a wire}\label{ch2:sec:motor}
\covers{M6.IQ2.01}

\begin{bigidea}
A wire carrying a current in an external magnetic field feels a force. The force is largest when
the wire is at \ang{90} to the field. There is no force when the wire is parallel to the field.
\end{bigidea}

\begin{keywords}
  \kw{Motor effect}{the force on a current-carrying wire in an external magnetic field.}
  \kw{External field}{a field from something else (such as a magnet), not from the wire.}
  \kw{Perpendicular}{at \ang{90} to something.}
  \kw{Parallel}{side by side and pointing the same way (or exactly the opposite way), like train
    rails. The angle between them is \ang{0} (or \ang{180}).}
  \kw{Right-hand palm rule}{fingers along $B$, thumb along $I$: the palm pushes along $F$.}
\end{keywords}

Motors, loudspeakers and electric toothbrushes all use the \term{motor effect}. A wire carrying a
current in an \term{external magnetic field} feels a force. ``External'' means the field comes
from something else, such as a magnet.

Why? A current is a stream of moving charges. A moving charge in a magnetic field feels a force,
$F = qv_\perp B$ (\cref{ch1:chap}). The charges cannot leave the wire, so the whole wire is pushed.

The force grows with the length of wire in the field, the current and the field strength. Only the
part of the current that is \term{perpendicular} (at \ang{90}) to the field lines counts
(\cref{ch2:fig:cases}). If the wire is \term{parallel} to the field (it points along the field
lines, at \ang{0} or \ang{180}), there is no force.

\begin{figure}
  \centering
  \begin{tikzpicture}[x=1cm, y=1cm]
    \foreach \dx in {0, 4.7, 9.4} {
      \begin{scope}[xshift=\dx cm]
        \foreach \y in {-1.0,-0.33,0.33,1.0} {\draw[field, hscfield!70] (-1.7,\y) -- (1.7,\y);}
        \node[diagramlabel, text=hscfield] at (1.9,1.0) {$B$};
      \end{scope}
    }
    % (a) perpendicular
    \draw[chTwowire] (0,-1.3) -- (0,1.3);
    \draw[chTwocurrent] (0,-0.25) -- (0,0.7) node[left, inner sep=3pt] {$I$};
    \pic at (0.55,-0.67) {chTwoFin};
    \node[diagramlabel, text=hscforce, anchor=west] at (0.78,-0.67) {$F$};
    \node[diagramlabel, align=center] at (0,-2.15) {(a) $\theta = \ang{90}$:\\ largest force, $F = lIB$};
    % (b) parallel
    \begin{scope}[xshift=4.7cm]
      \draw[chTwowire] (-1.4,0) -- (1.4,0);
      \draw[chTwocurrent] (-0.6,0) -- (0.5,0) node[above, inner sep=3pt] {$I$};
      \node[diagramlabel, align=center] at (0,-2.15) {(b) $\theta = \ang{0}$:\\ no force};
    \end{scope}
    % (c) at an angle
    \begin{scope}[xshift=9.4cm]
      \coordinate (A) at (-0.95,-0.8); \coordinate (B) at (0.95,0.8);
      \draw[chTwowire] ($(A)!-0.12!(B)$) -- ($(A)!1.12!(B)$);
      \draw[chTwocurrent] (A) -- (B) node[pos=0.62, below right, inner sep=1pt, fill=white] {$I$};
      \draw[->, line width=1pt, hscacc] (A) -- (0.95,-0.8) node[pos=0.5, below=9pt, inner sep=1pt, text=hscacc] {$I_\parallel = I\cos\theta$};
      \draw[->, line width=1pt, hscacc] (0.95,-0.8) -- (B) node[pos=0.45, right, inner sep=2pt, text=hscacc, fill=white] {$I_\perp = I\sin\theta$};
      \draw[hscink, line width=0.6pt] ($(A)+(0.55,0)$) arc (0:40.1:0.55);
      \node[diagramlabel, fill=white, inner sep=1pt] at ($(A)+(0.78,0.26)$) {$\theta$};
      \node[diagramlabel, align=center] at (0,-2.15) {(c) angle $\theta$: only $I_\perp$\\ counts, $F = lIB\sin\theta$};
    \end{scope}
  \end{tikzpicture}
  \caption{A wire in a uniform field that points to the right. In (a) the force is into the page.}\label{ch2:fig:cases}
\end{figure}

\begin{toolkit}{Splitting a vector with sin and cos}
\textbf{sin and cos.} Take a right-angled triangle and one of its other angles, $\theta$. The
longest side is the \emph{hypotenuse}. Then $\sin\theta = \text{opposite} \div
\text{hypotenuse}$ and $\cos\theta = \text{adjacent} \div \text{hypotenuse}$.\\
\textbf{Splitting a current.} Draw the current $I$ as the hypotenuse (\cref{ch2:fig:cases}(c)).
The side along the field is next to $\theta$: $I_\parallel = I\cos\theta$. The side across the
field is opposite $\theta$: $I_\perp = I\sin\theta$.\\
\textbf{Which angle?} $\theta$ is the angle between the wire and the field lines. Sometimes the
angle you are given is measured from a line at \ang{90} to the field. Then
$\theta = \ang{90} - (\text{given angle})$.\\
\textbf{Values to know:} $\sin\ang{0} = 0$, $\sin\ang{30} = 0.5$, $\sin\ang{90} = 1$. Set your
calculator to degrees (D or DEG on the screen).
\end{toolkit}

\begin{sidebyside}[0.44]
  \centering
  \begin{tikzpicture}
    \begin{axis}[hsc, width=0.93\linewidth, height=0.56\linewidth,
        xmin=0, xmax=180, ymin=0, ymax=1.15, xtick={0,45,90,135,180},
        ytick={0,0.5,1}, xlabel={$\theta$ (degrees)}, ylabel={$F \div F_\mathrm{max}$}]
      \addplot[hscseriesA, line width=1.2pt, domain=0:180, samples=91, mark=none] {sin(x)};
      \addplot[hscseriesA, only marks, mark=*, mark size=2pt] coordinates {(0,0) (90,1) (180,0)};
    \end{axis}
  \end{tikzpicture}
  \captionof{figure}{The force follows $\sin\theta$. $F_\mathrm{max} = lIB$ is the largest force,
    at $\theta = \ang{90}$.}\label{ch2:fig:sine}
\nextside
  \Cref{ch2:fig:sine} shows how the force changes with the angle.
  \begin{itemize}
    \item \textbf{Largest force:} wire at \ang{90} to the field ($\sin\ang{90} = 1$): $F = lIB$.
    \item \textbf{No force:} the wire is parallel to the field ($\theta = \ang{0}$ or \ang{180}).
          There is also no force when there is no current, or when the wire is outside the field.
    \item \textbf{Special case:} if $B$ points into or out of the page, any wire lying in the page
          is at \ang{90} to it. So use $\sin\theta = 1$, whatever angle the wire makes on the page.
  \end{itemize}
\end{sidebyside}

\textbf{Direction of the force.} The force is at \ang{90} to both the current and the field.

\begin{sidebyside}[0.36]
  \centering
  \begin{tikzpicture}[x=1cm, y=1cm]
    \coordinate (O) at (0,0);
    \coordinate (X) at (1.6,0); \coordinate (Y) at (0,1.6); \coordinate (Z) at (-0.85,-0.7);
    % the corner of a box
    \draw[guide] (X) -- ++(0,1.6) -- (Y);
    \draw[guide] (X) -- ++(-0.85,-0.7) -- (Z);
    \draw[guide] (Y) -- ++(-0.85,-0.7) -- (Z);
    % right-angle marks
    \draw[hscink, line width=0.5pt] (0.25,0) -- (0.25,0.25) -- (0,0.25);
    \draw[hscink, line width=0.5pt] (0.25,0) -- (0.117,-0.11) -- (-0.133,-0.11);
    \draw[hscink, line width=0.5pt] (0,0.25) -- (-0.133,0.14) -- (-0.133,-0.11);
    % the three directions
    \draw[chTwocurrent] (O) -- (X) node[right, inner sep=2pt] {$I$};
    \draw[->, line width=1.3pt, hscfield] (O) -- (Y) node[above, inner sep=2pt] {$B$};
    \draw[force] (O) -- (Z) node[below left, inner sep=1pt, text=hscforce] {$F$};
  \end{tikzpicture}
  \captionof{figure}{$F$ is at \ang{90} to both $I$ and $B$.}\label{ch2:fig:axes}
\nextside
  Picture the corner of a box (\cref{ch2:fig:axes}). $I$, $B$ and $F$ lie along its three edges,
  so each one is at \ang{90} to the other two. In the picture, $I$ points to the right and $B$
  points up the page. So $F$ comes out of the page, towards you.
\end{sidebyside}

To find which way $F$ points, use the \term{right-hand palm rule} (\cref{ch2:fig:palm}). It is the
same rule as for a moving positive charge in \cref{ch1:chap}. The thumb goes along the
conventional current, not the electron flow.

\begin{sidebyside}[0.50]
  \centering
  \begin{tikzpicture}[x=1cm, y=1cm]
    \tikzset{hand/.style={draw=hscink, line width=0.8pt, fill=hscbrand!7}}
    % thumb, fingers, wrist, then palm on top
    \draw[hand, rounded corners=6pt] (1.45,0.32) rectangle (2.95,0.82);
    \draw[hand, rounded corners=5pt] (0.05,1.6) rectangle (0.43,2.8);
    \draw[hand, rounded corners=5pt] (0.47,1.6) rectangle (0.86,3.2);
    \draw[hand, rounded corners=5pt] (0.90,1.6) rectangle (1.29,3.35);
    \draw[hand, rounded corners=5pt] (1.33,1.6) rectangle (1.72,3.12);
    \draw[hand] (0.3,-0.75) rectangle (1.45,0.2);
    \draw[hand, rounded corners=9pt] (0,0) rectangle (1.77,1.9);
    \pic at (0.88,0.95) {chTwoFout};
    % arrows and labels
    \draw[->, line width=1.3pt, hscfield] (-0.45,1.4) -- (-0.45,3.3);
    \node[diagramlabel, text=hscfield, align=right, anchor=east] at (-0.55,2.4) {fingers:\\ field $B$};
    \draw[chTwocurrent] (3.1,0.57) -- (4.1,0.57);
    \node[diagramlabel, align=center, anchor=north] at (3.6,0.4) {thumb:\\ current $I$};
    \node[diagramlabel, text=hscforce, align=left, anchor=west] at (2.0,1.75) {palm pushes:\\ force $F$ out\\ of the page};
    \draw[hscforce, line width=0.5pt] (1.97,1.6) -- (1.1,1.05);
  \end{tikzpicture}
  \captionof{figure}{The right-hand palm rule (palm facing you).}\label{ch2:fig:palm}
\nextside
  \begin{enumerate}
    \item Hold your right hand flat.
    \item \textbf{Fingers} along the field $B$ (N to S).
    \item \textbf{Thumb} along the conventional current $I$. Electron flow? Point the thumb the
          opposite way.
    \item The \textbf{palm} faces the way the force $F$ pushes.
  \end{enumerate}
  Reverse the current \emph{or} the field, and the force reverses. Reverse both, and it stays the same.
\end{sidebyside}

\begin{formulacard}{Force on a current-carrying wire}
  \cardformula{F = lI_\perp B = lIB\sin\theta}
  \begin{symbols}
    \sym{$F$}{size of the force on the wire}{\si{N}}
    \sym{$l$}{length of wire inside the field}{\si{m}}
    \sym{$I$}{current in the wire; $I_\perp = I\sin\theta$ is the part at \ang{90} to $B$}{\si{A}}
    \sym{$B$}{magnetic field strength of the external field}{\si{T} ($= \si{N.A^{-1}.m^{-1}}$)}
    \sym{$\theta$}{angle between the wire (the current) and the field}{degree (\si{\degree})}
  \end{symbols}
  \inwords{a longer wire, a bigger current or a stronger field gives a bigger force.}
  \whentouse{a straight wire with a current in a uniform external magnetic field.}
  \whennottouse{to find the direction (use the palm rule); for a curved wire or a non-uniform field
    (split it into short straight pieces); or for one moving charge ($F = qv_\perp B$).}
  \rearranged{$B = \dfrac{F}{lI\sin\theta}$ \qquad $I = \dfrac{F}{lB\sin\theta}$ \qquad
    $l = \dfrac{F}{IB\sin\theta}$ \qquad $\sin\theta = \dfrac{F}{lIB}$}
  \tensecond{$l = \qty{0.20}{m}$, $I = \qty{3.0}{A}$, $B = \qty{0.50}{T}$, $\theta = \ang{30}$:
    $F = 0.20 \times 3.0 \times 0.50 \times \sin\ang{30} = \qty{0.15}{N}$.}
  \trap{$\theta$ is the angle between the wire and the field lines. Sometimes the given angle is
    measured from a line at \ang{90} to the field (for example from the horizontal when $B$ is
    vertical). Then use $\theta = \ang{90} - (\text{given angle})$. Change cm to m and mA to A.}
\end{formulacard}

\begin{workedexample}{A wire in a field into the page}\label{ch2:we:basic}
A straight wire \qty{8.0}{cm} long lies in the page and carries \qty{2.5}{A} to the right. A
uniform field of \qty{45}{mT} points into the page. Find the size and direction of the force.
\given{$l = \qty{8.0}{cm} = \qty{0.080}{m}$, \ $I = \qty{2.5}{A}$, \ $B = \qty{45}{mT} = \qty{0.045}{T}$.
  The wire is in the page and $B$ is into the page, so $\theta = \ang{90}$.}
\find{the size and direction of the force $F$}
\formula{$F = lIB\sin\theta$, then the right-hand palm rule}
\substitute{$F = \qty{0.080}{m} \times \qty{2.5}{A} \times \qty{0.045}{T} \times \sin\ang{90}
  = \qty{9.0e-3}{N}$}
\answer{\final{$F = \qty{9.0e-3}{N}$ up the page}. (Hold the page upright in front of you. Point
  your fingers into the page, away from you, and your thumb to the right. Your palm faces up the
  page.)}
\checkstep{units: $\unit{m} \times \unit{A} \times \unit{N.A^{-1}.m^{-1}} = \unit{N}$ \cmark.
  $\theta = \ang{90}$, so this is the largest force this wire can feel. \qty{9.0e-3}{N} is about
  the weight of a \qty{0.9}{g} mass ($0.0090 \div 9.8 \approx \qty{9.2e-4}{kg}$). This is a
  small push, as expected for \qty{8}{cm} of wire carrying \qty{2.5}{A}.}
\end{workedexample}

\begin{workedexample}{The angle is measured from the horizontal}\label{ch2:we:angle}
\begin{sidebyside}[0.54]
A wire \qty{0.30}{m} long carries \qty{3.2}{A} in a uniform field of \qty{0.12}{T}. The field
points up the page. The wire makes \ang{25} with the horizontal, and the current flows to the
right and slightly up (\cref{ch2:fig:weangle}). Find the size and direction of the force on the
wire.
\nextside
  \centering
  \begin{tikzpicture}[x=1cm, y=1cm]
    \foreach \x in {0.45,1.45,2.45} {\draw[field] (\x,-0.45) -- (\x,2.0);}
    \node[diagramlabel, text=hscfield, anchor=south] at (1.45,2.05) {$B$};
    \draw[guide] (-0.3,0) -- (3.2,0);
    \draw[chTwowire] (0,0) -- (25:3.1);
    \draw[chTwocurrent] (25:1.75) -- (25:2.6);
    \node[diagramlabel, fill=white, inner sep=1pt] at ($(25:2.2)+(-0.1,0.32)$) {$I$};
    \draw[hscink, line width=0.6pt] (0.8,0) arc (0:25:0.8);
    \node[diagramlabel, fill=white, inner sep=1pt] at (11:1.22) {\ang{25}};
  \end{tikzpicture}
  \captionof{figure}{The wire makes \ang{25} with the dashed horizontal line.}\label{ch2:fig:weangle}
\end{sidebyside}
\given{$l = \qty{0.30}{m}$, \ $I = \qty{3.2}{A}$, \ $B = \qty{0.12}{T}$. The \ang{25} is measured
  from the dashed horizontal line. The field points up the page, so this line is at \ang{90} to
  the field. The wire is \ang{25} away from that line. So the angle between the wire and $B$ is
  $\theta = \ang{90} - \ang{25} = \ang{65}$.}
\find{the size and direction of the force $F$}
\formula{$F = lIB\sin\theta$, then the right-hand palm rule}
\substitute{$F = \qty{0.30}{m} \times \qty{3.2}{A} \times \qty{0.12}{T} \times \sin\ang{65}
  = \qty{0.1152}{N} \times 0.906 = \qty{0.104}{N}$}
\answer{\final{$F = \qty{0.10}{N}$ out of the page} (fingers up the page, thumb to the right and
  slightly up: the palm faces out of the page)}
\checkstep{it is less than the largest possible force, $lIB = \qty{0.115}{N}$ \cmark. Using
  $\sin\ang{25}$ by mistake gives \qty{0.049}{N}.}
\end{workedexample}

\begin{mistakes}
\begin{itemize}
  \item \xmark\ Putting the given angle straight into $\sin\theta$ when it is measured from a line
        at \ang{90} to the field. \cmark\ $\theta$ is between the wire and the field: use
        $\ang{90} - (\text{given angle})$.
  \item \xmark\ Thumb along the electron flow. \cmark\ Thumb along the conventional current.
  \item \xmark\ Force along the field. \cmark\ The force is at \ang{90} to both $B$ and $I$.
  \item \xmark\ Putting $l$ in cm or $I$ in mA. \cmark\ Change to m and A first.
\end{itemize}
\end{mistakes}

\begin{examtip}
\begin{itemize}
  \item ``Define the motor effect'' (1 mark): the answer must say \emph{external magnetic field}.
  \item ``Explain why there is no force'' (2 marks): (1) the current is parallel to the field, so
        $\theta = \ang{0}$; (2) $F = lIB\sin\theta$ and $\sin\ang{0} = 0$, so $F = 0$.
  \item ``Which way is the force (or the field)?'' (1 mark): the question often gives the electron
        flow. Reverse it first. Then use the palm rule. To find $B$, put your thumb along $I$ and
        your palm along $F$, and read your fingers.
  \item A force answer needs the size \emph{and} the direction (``into the page'', ``to the
        left'', never just ``sideways'').
\end{itemize}
\end{examtip}

\begin{quickcheck}
  \item Which angle between wire and field gives (a) the largest force, (b) no force?
  \item A \qty{12}{cm} wire carries \qty{4.0}{A} at \ang{90} to a \qty{0.30}{T} field. Find the
        force on it.\label{ch2:qc:wire}
  \item The wire in \cref{ch2:qc:wire} is turned to \ang{30} to the field. Find the new force.
  \item A field points to the left. The electrons in a wire flow out of the page. Which way is
        the force on the wire?
  \item A \qty{5.0}{cm} wire carrying \qty{2.0}{A} at \ang{90} to a field feels \qty{4.0e-3}{N}.
        Find the field strength.
\end{quickcheck}

%% ==================================================================
\section{Forces between parallel wires}\label{ch2:sec:wires}
\covers{M6.IQ2.03}

\begin{bigidea}
Each of two parallel wires sits in the field of the other, so each feels a motor-effect force.
Same-direction currents attract; opposite currents repel. The two forces are always equal and
opposite, and they were used to define the ampere.
\end{bigidea}

\begin{keywords}
  \kw{Attract / repel}{pull together / push apart.}
  \kw{Force per unit length ($F/l$)}{the force on each metre of wire. Unit: \si{N.m^{-1}}.}
  \kw{Permeability constant ($\mu_0$)}{links a current to the field it makes in a vacuum.}
  \kw{Newton's third law}{A pushes on B, so B pushes on A: same size, opposite way.}
  \kw{Ampere (A)}{the unit of current. \qty{1}{A} means one coulomb (C) of charge passes a point
    every second.}
\end{keywords}

Cables that carry big currents side by side push or pull on each other. Jump-start leads twitch
when the car starts. In power stations, the thick metal bars that carry the current must be
bolted down firmly.

Wire 1 makes a magnetic field. Wire 2 is a current in that field, so it feels a motor-effect
force. The same is true the other way round. Find the direction in three steps
(\cref{ch2:fig:parallel}):
\begin{enumerate}
  \item grip rule: find the field of wire 1 at wire 2;
  \item palm rule: find the force on wire 2;
  \item Newton's third law: the force on wire 1 is the same size, in the opposite direction.
\end{enumerate}

\begin{figure}
  \centering
  \begin{tikzpicture}[x=1cm, y=1cm]
    \foreach \dx in {0, 7.2} {
      \begin{scope}[xshift=\dx cm]
        \draw[guide] (-0.55,0) -- (3.55,0);
        \draw[chTwoarc] (-24:3) arc (-24:24:3);
        \draw[->, line width=1.2pt, hscfield] (3,0.2) -- (3,1.0) node[right, inner sep=2pt] {$B_1$};
      \end{scope}
    }
    % (a) same direction
    \draw[chTwoarc] (156:3) ++(3,0) arc (156:204:3);
    \draw[->, line width=1.2pt, hscfield] (0,-0.2) -- (0,-1.0) node[left, inner sep=2pt] {$B_2$};
    \pic at (0,0) {chTwowireout}; \pic at (3,0) {chTwowireout};
    \draw[force] (0.25,0) -- (1.2,0) node[midway, above, inner sep=2pt, text=hscforce] {$F$};
    \draw[force] (2.75,0) -- (1.8,0) node[midway, above, inner sep=2pt, text=hscforce] {$F$};
    \node[diagramlabel, align=right, anchor=east] at (-0.3,0.4) {wire 1\\ \qty{2.0}{A}};
    \node[diagramlabel, align=left, anchor=west] at (3.3,-0.4) {wire 2\\ \qty{6.0}{A}};
    \node[diagramlabel] at (1.5,-1.75) {(a) same direction: \textbf{attract}};
    % (b) opposite directions
    \begin{scope}[xshift=7.2cm]
      \draw[chTwoarc] (204:3) ++(3,0) arc (204:156:3);
      \draw[->, line width=1.2pt, hscfield] (0,0.2) -- (0,1.0) node[left, inner sep=2pt] {$B_2$};
      \pic at (0,0) {chTwowireout}; \pic at (3,0) {chTwowirein};
      \draw[force] (-0.25,0) -- (-1.25,0) node[above, inner sep=2pt, text=hscforce] {$F$};
      \draw[force] (3.25,0) -- (4.25,0) node[above, inner sep=2pt, text=hscforce] {$F$};
      \node[diagramlabel, anchor=north east] at (-0.2,-0.25) {wire 1};
      \node[diagramlabel, anchor=north west] at (3.2,-0.25) {wire 2};
      \node[diagramlabel] at (1.5,-1.75) {(b) opposite directions: \textbf{repel}};
    \end{scope}
  \end{tikzpicture}
  \caption{Two wires seen end-on. $B_1$ is the field made by wire 1, at wire 2. $B_2$ is the field
    made by wire 2, at wire 1. Each curved line is part of the circular field line round the other
    wire. In (a) the currents differ, but the two forces are the same size.}\label{ch2:fig:parallel}
\end{figure}

In \cref{ch2:fig:parallel}(a) both currents come out of the page. Wire 1's field lines go
anticlockwise, so at wire 2 the field points up. Palm rule on wire 2: thumb out of the page,
fingers up, palm facing left, towards wire 1. The wires attract.

So currents in the same direction (\textbf{like currents}) attract. This is the opposite of
electric charges, where like charges repel.

You can also picture it with the fields. Between two same-direction currents, the two fields
point opposite ways, so they partly cancel. The field there is weak. The wires move towards this
weaker region. (This picture is only a memory aid. Always find the direction with the grip rule
and the palm rule.)

\textbf{How big is the force?} Wire 2 sits in wire 1's field, $B_1 = \mu_0 I_1/(2\pi r)$. A bigger
$I_1$ makes a stronger field. A bigger $I_2$ means more current for that field to push. A bigger
gap $r$ makes the field weaker. So the force grows with $I_1$ and $I_2$, and gets smaller as $r$
grows.

\begin{formulacard}{Force between two parallel wires}
  \cardformula{\frac{F}{l} = \frac{\mu_0}{2\pi}\frac{I_1 I_2}{r}}
  \begin{symbols}
    \sym{$F$}{force on \emph{each} wire}{\si{N}}
    \sym{$l$}{length over which the wires run side by side}{\si{m}}
    \sym{$\mu_0$}{magnetic permeability constant; $\mu_0/2\pi = 2 \times 10^{-7}$}{\si{N.A^{-2}}}
    \sym{$I_1$, $I_2$}{the currents in the two wires}{\si{A}}
    \sym{$r$}{distance between the wires (centre to centre)}{\si{m}}
  \end{symbols}
  \inwords{the force on each metre grows with both currents and falls as the wires move apart.}
  \whentouse{two long, straight, parallel wires (much longer than the gap $r$).}
  \whennottouse{to decide attract or repel (use the rules), or for wires that are not parallel.}
  \rearranged{$F = \dfrac{\mu_0 I_1 I_2\, l}{2\pi r}$ \qquad $r = \dfrac{\mu_0 I_1 I_2\, l}{2\pi F}$
    \qquad $I_2 = \dfrac{2\pi r F}{\mu_0 I_1 l}$}
  \tensecond{$I_1 = \qty{3.0}{A}$, $I_2 = \qty{4.0}{A}$, $r = \qty{0.020}{m}$:
    $\dfrac{F}{l} = \dfrac{2 \times 10^{-7} \times 3.0 \times 4.0}{0.020} = \qty{1.2e-4}{N.m^{-1}}$.}
  \derive{(HSC questions ask for this) wire 1 makes a field $B_1 = \dfrac{\mu_0 I_1}{2\pi r}$ at
    wire 2. $B_1$'s field lines are circles round wire 1. Wire 2 runs straight, parallel to
    wire 1, so it is at \ang{90} to these circles. So $\sin\theta = 1$.
    \begin{align*}
      F &= l I_2 B_1 && \text{motor effect on wire 2, with $\sin\theta = 1$}\\
        &= l I_2 \frac{\mu_0 I_1}{2\pi r} && \text{replace $B_1$}\\
        &= \frac{\mu_0}{2\pi}\frac{I_1 I_2\, l}{r} && \text{put the constants $\mu_0$ and $2\pi$ together}\\
      \frac{F}{l} &= \frac{\mu_0}{2\pi}\frac{I_1 I_2}{r} && \text{divide both sides by $l$}
    \end{align*}}
  \trap{the formula gives $F/l$ in \si{N.m^{-1}}. Multiply by $l$ to get the force. It is $r$, not $r^2$.}
\end{formulacard}

\textbf{Newton's third law for wires.} \term{Newton's third law} says that if A pushes on B, then B
pushes on A equally hard in the opposite direction. The formula uses the product
$I_1 \times I_2$. Swap the two wires, and the product stays the same ($2.0 \times 6.0 =
6.0 \times 2.0$). So both wires feel the \emph{same size} of force, even when the currents
differ. In
\cref{ch2:fig:parallel}(a) the \qty{6.0}{A} wire does \emph{not} pull harder than the \qty{2.0}{A}
wire. The two forces are an \term{action--reaction pair}: two equal and opposite forces that two
objects put on each other.

\textbf{The SI ampere.} The \term{SI} is the International System of Units: the set of units that
scientists everywhere agree to use. From 1948 to 2019, the official definition of the
\term{ampere} (the SI unit of current) used this force. Learn the definition in parts:
\begin{itemize}
  \item Take two wires. They are straight and parallel.
  \item They are infinitely long and very thin (\emph{negligible circular cross-section}).
  \item They are \qty{1}{m} apart in a vacuum (empty space, with no air).
  \item The same steady current flows in each wire.
\end{itemize}
If the force between the wires is exactly $2 \times 10^{-7}$ N on each metre of length, that
current is \qty{1}{A}.

\textbf{Exam wording (learn it):} \emph{one ampere is the constant current that, flowing in each
of two straight, parallel wires of infinite length and negligible circular cross-section, placed
\qty{1}{m} apart in a vacuum, produces \textbf{between} these wires a force of
$2 \times 10^{-7}$ newton per metre of length.}

Check with the formula: $\dfrac{F}{l} = 2 \times 10^{-7} \times \dfrac{1 \times 1}{1} =
\qty{2e-7}{N.m^{-1}}$ \cmark.

Once the ampere was defined, the coulomb was defined from it. One coulomb is the charge carried
past a point by a current of \qty{1}{A} in \qty{1}{s}.

\textbf{The link to Newton's third law.} The definition names one force \emph{between} the wires.
It does not say which wire. This works because of Newton's third law. Each wire pulls (or pushes)
on the other with a force of the same size, in the opposite direction. So one value,
\qty{2e-7}{N} per metre, is true for both wires and can define the unit. (Since 2019 the SI fixes the size of the
electron's charge instead. The force rule still holds, and HSC questions use the wire definition.)

\begin{workedexample}{Jump-start leads}\label{ch2:we:leads}
Two jump-start leads carry current to and from a car battery. For \qty{0.80}{m} they lie straight
and side by side, \qty{1.5}{cm} apart (centre to centre). Each carries \qty{120}{A}, in opposite
directions. Find the force on each lead. Do they attract or repel?
\given{$I_1 = I_2 = \qty{120}{A}$, \ $r = \qty{1.5}{cm} = \qty{0.015}{m}$, \ $l = \qty{0.80}{m}$,
  \ currents opposite}
\find{the force $F$ on each lead; attract or repel}
\formula{$\dfrac{F}{l} = \dfrac{\mu_0}{2\pi}\dfrac{I_1 I_2}{r}$, \ then $F = \dfrac{F}{l} \times l$}
\substitute{\eqlines{\frac{F}{l} &= 2 \times 10^{-7}\,\unit{N.A^{-2}} \times
  \frac{\qty{120}{A} \times \qty{120}{A}}{\qty{0.015}{m}} = \qty{0.192}{N.m^{-1}}\\
  F &= \qty{0.192}{N.m^{-1}} \times \qty{0.80}{m} = \qty{0.154}{N}}}
\answer{\final{$F = \qty{0.15}{N}$ on each lead}; opposite currents, so the leads \textbf{repel}}
\checkstep{units: $\unit{N.A^{-2}} \times \unit{A^2} \div \unit{m} \times \unit{m} = \unit{N}$
  \cmark. \qty{0.15}{N} is about the weight of a \qty{15}{g} mass ($0.15 \div 9.8 \approx
  \qty{0.015}{kg}$): enough to make loose leads twitch.}
\end{workedexample}

\begin{workedexample}{Two tracks on a circuit board}\label{ch2:we:overlap}
\begin{sidebyside}[0.56]
On a circuit board, current flows along thin, straight copper strips called tracks. Track X is
\qty{15.0}{cm} long and carries \qty{600}{mA} up the page. Track Y is \qty{6.0}{cm} long and
carries \qty{1.5}{A} down the page. Y is parallel to X and \qty{2.0}{mm} to its right. Y is level
with the middle part of X (\cref{ch2:fig:tracks}). (a) Find the force on Y. (b) Find the force
on X.
\nextside
  \centering
  \begin{tikzpicture}[x=1cm, y=1cm]
    \draw[guide] (-0.35,0.9) -- (2.1,0.9);
    \draw[guide] (-0.35,2.1) -- (2.1,2.1);
    \draw[chTwowire] (0,0) -- (0,3.0);
    \draw[chTwocurrent] (0,0.15) -- (0,0.8);
    \node[diagramlabel, anchor=east] at (-0.12,0.45) {\qty{600}{mA}};
    \node[diagramlabel, anchor=south] at (0,3.05) {X (\qty{15.0}{cm})};
    \draw[chTwowire] (1.0,0.9) -- (1.0,2.1);
    \draw[chTwocurrent] (1.0,1.95) -- (1.0,1.05);
    \node[diagramlabel, anchor=north] at (1.0,0.84) {\qty{1.5}{A}};
    \node[diagramlabel, anchor=south] at (1.0,2.15) {Y};
    \draw[<->, hscmuted, line width=0.5pt] (2.0,0.9) -- (2.0,2.1);
    \node[diagramlabel, anchor=west, align=left] at (2.05,1.5) {$l = \qty{6.0}{cm}$\\ (side by side)};
    \draw[<->, hscmuted, line width=0.5pt] (0.06,1.65) -- (0.94,1.65);
    \node[diagramlabel, anchor=south] at (0.5,1.65) {$r$};
  \end{tikzpicture}
  \captionof{figure}{Only \qty{6.0}{cm} of the two tracks is side by side. The gap
    $r = \qty{2.0}{mm}$ is not to scale.}\label{ch2:fig:tracks}
\end{sidebyside}
\given{$I_\mathrm{X} = \qty{600}{mA} = \qty{0.600}{A}$, \ $I_\mathrm{Y} = \qty{1.5}{A}$, \
  $r = \qty{2.0}{mm} = \qty{2.0e-3}{m}$. The tracks are side by side only along Y, so
  $l = \qty{6.0}{cm} = \qty{0.060}{m}$.}
\find{the size and direction of the force on each track}
\formula{$F = \dfrac{\mu_0}{2\pi}\dfrac{I_\mathrm{X} I_\mathrm{Y}\, l}{r}$ \ (the sheet formula
  multiplied by $l$)}
\substitute{$F = 2 \times 10^{-7}\,\unit{N.A^{-2}} \times
  \dfrac{\qty{0.600}{A} \times \qty{1.5}{A} \times \qty{0.060}{m}}{\qty{2.0e-3}{m}} = \qty{5.4e-6}{N}$}
\answer{(a) \final{$\qty{5.4e-6}{N}$ to the right} (opposite currents repel, so Y is pushed away
  from X). (b) \final{$\qty{5.4e-6}{N}$ to the left}: the same size (Newton's third law), even
  though Y's current is bigger.}
\checkstep{using \qty{15.0}{cm} would give \qty{1.4e-5}{N}, which is wrong: only \qty{6.0}{cm} of
  X lies beside Y. The tracks are long compared with the gap ($\qty{6.0}{cm} \div \qty{2.0}{mm} =
  30$), so the formula works \cmark. Small currents give a tiny force \cmark.}
\end{workedexample}

\begin{mistakes}
\begin{itemize}
  \item \xmark\ ``Like currents repel.'' \cmark\ Same-direction currents \emph{attract}.
  \item \xmark\ Giving $F/l$ as the force. \cmark\ Multiply by $l$, the length where the wires are
        side by side.
  \item \xmark\ ``The bigger current feels the bigger force.'' \cmark\ Both forces are the same size.
  \item \xmark\ Using $r^2$, or $r$ in cm. \cmark\ $r$ in metres, centre to centre.
\end{itemize}
\end{mistakes}

\begin{examtip}[breakable]
\begin{itemize}
  \item \textbf{Explain why the wires attract} (3 marks): (1) the field of wire 1 at wire 2, with
        its direction (grip rule); (2) the motor-effect force on wire 2, with its direction (palm
        rule: towards wire 1); (3) Newton's third law: wire 1 feels an equal and opposite force,
        so the wires attract.
  \item \textbf{Derive} $F/l$ (2 marks): (1) write $B_1 = \dfrac{\mu_0 I_1}{2\pi r}$ and
        $F = lI_2B_1$; (2) substitute and divide both sides by $l$. Keep the currents as two
        separate symbols, $I_1$ and $I_2$ (not $I^2$).
  \item \textbf{Define the ampere} (2 marks): a constant current of \qty{1}{A} in each of two
        straight, parallel, infinitely long wires of negligible cross-section, \qty{1}{m} apart
        in a vacuum, gives
        a force between them of $2 \times 10^{-7}$ N per metre of length.
  \item \textbf{Ampere and Newton's third law} (3 marks): (1) the definition names one force
        \emph{between} the wires; (2) by Newton's third law, each wire pushes or pulls on the other
        with an equal and opposite force, even if the currents differ (the product
        $I_1 \times I_2$ is the same for both wires); (3) so the same $2 \times 10^{-7}$ N per
        metre acts on each wire, and the definition does not depend on which wire you measure.
  \item \textbf{Define $\mu_0$} (1--2 marks): the magnetic permeability constant,
        $4\pi \times 10^{-7}$ \si{N.A^{-2}}. It links a current to the magnetic field it makes in
        a vacuum.
\end{itemize}
\end{examtip}

\begin{quickcheck}
  \item Two parallel wires carry currents in the same direction. Do they attract or repel?
        Explain using the motor effect.
  \item Two long parallel wires \qty{2.0}{cm} apart each carry \qty{5.0}{A}. Find $F/l$.\label{ch2:qc:pair}
  \item In \cref{ch2:qc:pair} the gap halves and one current doubles. By what factor does $F/l$ change?
  \item Wire A carries \qty{1.0}{A} and wire B carries \qty{4.0}{A}. A student says ``B pulls on A
        four times harder than A pulls on B.'' Explain why the student is wrong.
  \item State the (pre-2019) SI definition of the ampere. How does it depend on Newton's third law?
\end{quickcheck}

%% ==================================================================
\section{Key experiments}\label{ch2:sec:experiments}
\covers{M6.IQ2.01, M6.IQ2.02}

\begin{keywords}
  \kw{Electronic balance}{a digital scale; here it measures a tiny force.}
  \kw{Tare}{set the balance to read zero.}
  \kw{Validity}{does the experiment test what it claims to test?}
  \kw{Reliability}{do repeated readings agree?}
  \kw{Accuracy}{how close is the result to the true value?}
\end{keywords}

Both investigations use an \term{electronic balance} (a digital scale) to measure a tiny force.
To \term{tare} the balance means to set it to read zero.

\begin{refresher}{Weight and a balance reading}
A mass $m$ has weight $F = mg$. This is $F_\mathrm{net} = ma$ with $a = g$. It is not printed on
the formula sheet, so learn it.\\
A balance really measures the force on its pan. It divides this force by $g$ and shows the result
as a mass. $\Delta$ (the Greek letter delta) means ``change in''. So an extra force $F$ changes
the reading by $\Delta m = F/g$. Turned round, $F = \Delta m\, g$, with $\Delta m$ in kg.
\end{refresher}

\begin{toolkit}{Gradient of a straight-line graph}
Draw the \term{line of best fit}: the straight line that passes as close as possible to all the
points. Pick two points \emph{on the line}, far apart. Then
gradient $= \text{rise} \div \text{run} = (y_2 - y_1) \div (x_2 - x_1)$. Its unit is the $y$ unit
$\div$ the $x$ unit. If theory says $y = kx$, the gradient is $k$.
\end{toolkit}

\begin{experiment}{The force on a wire (current balance)}
\aim{To find how the force on a wire depends on the current, and when it is largest or zero.}
\apparatus{U-shaped magnet, electronic balance (reads to \qty{0.01}{g}), stiff straight wire,
  clamp stand, low-voltage DC supply (DC = direct current: it always flows the same way), ammeter
  (measures the current), \term{rheostat} (a variable resistor that sets the current), switch
  (\cref{ch2:fig:currentbalance}).
  \begin{figure}
  \centering
  \begin{tikzpicture}[x=1cm, y=1cm]
    % balance
    \draw[object, rounded corners=3pt] (0,0) rectangle (5.6,0.7);
    \draw[hscink, fill=white, line width=0.6pt] (0.25,0.15) rectangle (1.55,0.55);
    \node[font=\scriptsize] at (0.9,0.35) {\num{0.00} g};
    \draw[object] (1.4,0.7) rectangle (4.4,0.85);
    \node[diagramlabel, anchor=west] at (2.0,0.35) {electronic balance};
    % magnet
    \draw[draw=hscink, fill=hscmuted!25, line width=0.8pt] (1.7,0.85) -- (4.1,0.85) -- (4.1,2.6)
      -- (3.6,2.6) -- (3.6,1.25) -- (2.2,1.25) -- (2.2,2.6) -- (1.7,2.6) -- cycle;
    \draw[draw=hscink, fill=hscforce!20, line width=0.8pt] (2.2,1.5) rectangle (2.52,2.4);
    \draw[draw=hscink, fill=hscvel!20, line width=0.8pt] (3.28,1.5) rectangle (3.6,2.4);
    \node[font=\small\bfseries] at (2.36,1.95) {N};
    \node[font=\small\bfseries] at (3.44,1.95) {S};
    \foreach \y in {1.63,2.27} {\draw[field] (2.54,\y) -- (3.26,\y);}
    \pic at (2.9,1.95) {chTwowireout};
    % clamp stand holding the wire
    \draw[object] (-1.4,0) rectangle (-0.4,0.15);
    \draw[line width=2pt, hscmuted] (-0.9,0.15) -- (-0.9,3.2);
    \draw[line width=2pt, hscmuted] (-0.9,3.0) -- (2.9,3.0) -- (2.9,2.15);
    \node[diagramlabel, anchor=south west, align=left] at (-0.8,3.03) {the clamp holds the wire;\\
      the wire does not touch the magnet};
    % forces
    \draw[force] (3.1,2.2) -- (3.1,2.85) node[right, inner sep=2pt, text=hscforce] {$F$ on wire};
    \draw[force] (4.35,1.9) -- (4.35,1.2) node[right, inner sep=2pt, text=hscforce] {$F$ on magnet};
    % circuit: one series loop; the two leads go to the ends of the wire
    \begin{scope}[xshift=7.0cm, yshift=0.45cm]
      \draw[line width=0.8pt] (2.4,0.65) -- (2.4,2.4) -- (0,2.4) -- (0,0) -- (2.4,0) -- (2.4,0.15);
      \draw[line width=0.8pt] (2.4,0.15) -- (2.9,0.15) (2.4,0.65) -- (2.9,0.65);
      \fill (2.9,0.15) circle (1.3pt); \fill (2.9,0.65) circle (1.3pt);
      \node[diagramlabel, anchor=west, align=left] at (3.0,0.4) {to the two\\ ends of the wire};
      \node[draw, fill=white, rounded corners=2pt, font=\small] at (1.2,2.4) {DC supply};
      \node[draw, circle, fill=white, inner sep=1.5pt, font=\small] at (0,1.2) {A};
      \node[draw, fill=white, minimum width=3mm, minimum height=7mm, inner sep=0pt] (R) at (2.4,1.5) {};
      \draw[->, line width=0.6pt] ($(R)+(-0.3,-0.45)$) -- ($(R)+(0.3,0.45)$);
      \node[diagramlabel, anchor=west] at (2.7,1.5) {rheostat};
      \fill[white] (0.85,-0.1) rectangle (1.55,0.1);
      \draw[line width=0.8pt] (0.85,0) -- (1.5,0.3);
      \fill (0.85,0) circle (1.2pt); \fill (1.55,0) circle (1.2pt);
      \node[diagramlabel, anchor=north] at (1.2,-0.1) {switch};
    \end{scope}
  \end{tikzpicture}
  \caption{A current balance (front view; the wire is seen end-on, current out of the page). The
    supply, ammeter, rheostat, switch and wire form one series loop.}\label{ch2:fig:currentbalance}
  \end{figure}}
\method{\begin{enumerate}
  \item Clamp the wire through the gap at \ang{90} to the field. The wire must not touch the
        magnet. Connect it \emph{in series} (in one single loop) with the ammeter, rheostat,
        switch and supply. Tare the balance.
  \item Set $I = \qty{1.0}{A}$, record the change in reading $\Delta m$, and switch off. Repeat
        for 2.0, 3.0, 4.0 and \qty{5.0}{A}, three times each.
  \item Reverse the current. The reading now falls, because the force on the wire reverses. Then
        turn the magnet round (swap N and S): the force reverses again. Reverse both the current
        and the magnet: the force is the same as at first.
  \item To test the angle, use a short wire (\qty{3.0}{cm} long) that sits wholly inside the
        field. Its leads come straight down into the gap, so any force on them is sideways and
        does not change the reading. Keep $I = \qty{5.0}{A}$. Turn the magnet on the balance to
        change the angle $\theta$ between the wire and the field. Record $\Delta m$ at
        $\theta = \ang{0}$, \ang{30}, \ang{60} and \ang{90}.
\end{enumerate}
Why does the reading rise? The wire is pushed up (palm rule: fingers from N to S, thumb out of
the page, palm faces up). By Newton's third law, the magnet is pushed down with an equal force.
So the reading rises, and $F = \Delta m\, g$.

Quick qualitative demo: hang a loose strip of aluminium foil through the gap and connect its ends
to a low-voltage supply. The strip jumps when you switch on. Reverse the current or the magnet,
and it jumps the other way.}
\results{\begin{sidebyside}[0.34]
    \centering\small
    \begin{tabular}{@{}ccc@{}}
      \toprule
      $I$ (\si{A}) & $\Delta m$ (\si{g}) & $F$ ($10^{-3}$ \si{N})\\
      \midrule
      1.0 & 0.42 & 4.1\\
      2.0 & 0.88 & 8.6\\
      3.0 & 1.29 & 12.6\\
      4.0 & 1.75 & 17.2\\
      5.0 & 2.16 & 21.2\\
      \bottomrule
    \end{tabular}
  \nextside
    \centering
    \begin{tikzpicture}
      \begin{axis}[hsc, width=\linewidth, height=0.40\linewidth,
          xmin=0, xmax=5.5, ymin=0, ymax=24, xtick={0,1,...,5}, ytick={0,5,...,20},
          xlabel={$I$ (\si{A})}, ylabel={$F$ ($10^{-3}$ \si{N})}]
        \addplot[hscseriesA, only marks, mark=*, mark size=2pt]
          coordinates {(1,4.1) (2,8.6) (3,12.6) (4,17.2) (5,21.2)};
        \addplot[hscseriesB, dashed, line width=1pt, domain=0:5.4, samples=2, mark=none] {4.25*x};
      \end{axis}
    \end{tikzpicture}
    \captionof{figure}{Force against current, with the line of best fit (dashed).}\label{ch2:fig:fvsi}
  \end{sidebyside}
  These are sample class results, with $l = \qty{0.050}{m}$ of wire in the field. The heading
  ``$F$ ($10^{-3}$ N)'' means the numbers are in units of $10^{-3}$ N: 4.1 means
  $\qty{4.1e-3}{N}$. Example (row 1): $\Delta m = \qty{0.42}{g} = \qty{4.2e-4}{kg}$, so
  $F = \qty{4.2e-4}{kg} \times \qty{9.8}{m.s^{-2}} = \qty{4.1e-3}{N}$.

  Angle test (short wire, $I = \qty{5.0}{A}$):
  \begin{center}\small
  \begin{tabular}{@{}lcccc@{}}
    \toprule
    $\theta$ & \ang{0} & \ang{30} & \ang{60} & \ang{90}\\
    \midrule
    $\Delta m$ (\si{g}) & 0.00 & 0.66 & 1.12 & 1.29\\
    $\Delta m \div \qty{1.29}{g}$ & 0 & 0.51 & 0.87 & 1.00\\
    $\sin\theta$ & 0 & 0.50 & 0.87 & 1.00\\
    \bottomrule
  \end{tabular}
  \end{center}}
\conclusion{The points (\cref{ch2:fig:fvsi}) lie close to a straight line through the origin
  (0, 0), so $F \propto I$.
  Theory: $F = (lB)\,I$, so the gradient is $lB$. Read two points far apart on the line of best
  fit: (0, 0) and (\qty{5.0}{A}, \qty{21.25e-3}{N}). The gradient is
  $\qty{21.25e-3}{N} \div \qty{5.0}{A} = \qty{4.25e-3}{N.A^{-1}}$. So
  $B = \text{gradient} \div l = \qty{4.25e-3}{N.A^{-1}} \div \qty{0.050}{m} = \qty{0.085}{T}$.

  In the angle test, $\Delta m \div \qty{1.29}{g}$ matches $\sin\theta$. The force is largest at
  \ang{90} and zero when the wire is parallel to $B$. Reversing the current or the magnet reversed
  the force. All this agrees with $F = lIB\sin\theta$ and the palm rule.}
\significance{This is the motor effect behind every electric motor, loudspeaker and analogue
  meter. It also gives a way to measure $B$: $B = \text{gradient} \div l$.}
\validity{\begin{itemize}
  \item \textbf{Validity} (does it test what it claims?): change only $I$; keep $l$, $B$ and the
        angle fixed. Keep the wire in the middle of the gap, where the field is uniform. Keep the
        leads to the wire outside the gap. They carry the same current and would feel a force
        too.
  \item \textbf{Reliability} (do repeats agree?): repeat each reading three times and average.
  \item \textbf{Accuracy} (how close to the true value?): read $I$ from the ammeter, not from the
        supply dial. Switch on only briefly. A hot wire has more resistance, so the current drops
        while you read the balance.
  \item \textbf{Safety:} low voltage only; the wire gets hot; strong magnets pinch fingers.
\end{itemize}}
\end{experiment}

\begin{experiment}{Two parallel wires (required investigation)}
\aim{To show that parallel currents attract or repel, and to test $F/l = (\mu_0/2\pi)(I_1 I_2/r)$.}
\apparatus{two straight brass rods (\qty{0.50}{m}), electronic balance (reads to \qty{0.001}{g}),
  insulating blocks (they do not conduct), retort stand with an insulated clamp, flexible lead (a
  bendy wire), low-voltage high-current DC supply, ammeter, rheostat, switch, vernier calipers
  (they measure small gaps precisely) (\cref{ch2:fig:parallelrig}).
  \begin{figure}
  \centering
  \begin{tikzpicture}[x=1cm, y=1cm]
    % balance and blocks
    \draw[object, rounded corners=3pt] (0,0) rectangle (7,0.6);
    \draw[hscink, fill=white, line width=0.6pt] (0.3,0.12) rectangle (1.7,0.48);
    \node[font=\scriptsize] at (1.0,0.3) {\num{0.000} g};
    \node[diagramlabel, anchor=west] at (2.3,0.3) {electronic balance};
    \draw[object] (0.6,0.6) rectangle (6.4,0.72);
    \draw[draw=hscink, fill=hscmuted!30] (1.0,0.72) rectangle (1.4,1.18);
    \draw[draw=hscink, fill=hscmuted!30] (5.6,0.72) rectangle (6.0,1.18);
    % rods
    \draw[line width=3pt, hscmuted!80!black] (0.4,1.25) -- (6.6,1.25);
    \draw[line width=3pt, hscmuted!80!black] (0.4,2.05) -- (6.6,2.05);
    \draw[line width=0.8pt] (6.6,2.05) .. controls (7.2,2.05) and (7.2,1.25) .. (6.6,1.25);
    \node[diagramlabel, anchor=west, align=left] at (7.15,1.65) {flexible\\ link};
    % clamp from above
    \draw[line width=2pt, hscmuted] (3.5,3.1) -- (3.5,2.12);
    \draw[line width=2pt, hscmuted] (3.5,3.1) -- (8.9,3.1) -- (8.9,0.15);
    \draw[object] (8.4,0) rectangle (9.4,0.15);
    \node[diagramlabel, anchor=south] at (6.2,3.12) {retort stand, insulated clamp};
    % currents
    \draw[chTwocurrent] (1.9,2.3) -- (3.0,2.3) node[right, inner sep=2pt] {$I$};
    \draw[chTwocurrent] (4.9,1.0) -- (3.8,1.0) node[left, inner sep=2pt] {$I$};
    % dimensions
    \draw[<->, hscmuted, line width=0.5pt] (0.4,2.62) -- (6.6,2.62);
    \node[diagramlabel, fill=white, inner sep=1pt] at (5.1,2.62) {$l = \qty{0.50}{m}$};
    \draw[<->, hscmuted, line width=0.5pt] (4.3,1.33) -- (4.3,1.97);
    \node[diagramlabel, anchor=west] at (4.35,1.65) {$r$};
    \draw[force] (2.6,1.15) -- (2.6,0.78);
    \node[diagramlabel, text=hscforce, anchor=west] at (2.65,0.93) {$F$};
    % circuit (left side): supply on top, ammeter and rheostat on the left, switch below
    \draw[line width=0.8pt] (0.4,2.05) -- (-0.6,2.05) -- (-0.6,3.4) -- (-2.6,3.4) -- (-2.6,0.6)
      -- (-0.6,0.6) -- (-0.6,1.25) -- (0.4,1.25);
    \node[draw, fill=white, rounded corners=2pt, font=\small] at (-1.6,3.4) {DC supply};
    \node[draw, circle, fill=white, inner sep=1.5pt, font=\small] at (-2.6,2.65) {A};
    \node[draw, fill=white, minimum width=3mm, minimum height=7mm, inner sep=0pt] (R2) at (-2.6,1.5) {};
    \draw[->, line width=0.6pt] ($(R2)+(-0.3,-0.45)$) -- ($(R2)+(0.3,0.45)$);
    \node[diagramlabel, anchor=east] at (-2.95,1.5) {rheostat};
    \fill[white] (-1.95,0.5) rectangle (-1.25,0.7);
    \draw[line width=0.8pt] (-1.95,0.6) -- (-1.3,0.9);
    \fill (-1.95,0.6) circle (1.2pt); \fill (-1.25,0.6) circle (1.2pt);
    \node[diagramlabel, anchor=north] at (-1.6,0.5) {switch};
  \end{tikzpicture}
  \caption{Two rods in series carry opposite currents (gap not to scale).}\label{ch2:fig:parallelrig}
  \end{figure}}
\method{\begin{enumerate}
  \item Rest the lower rod on the insulating blocks on the balance. Clamp the upper rod directly
        above it and parallel to it. Measure the gap between the rod surfaces at both ends with
        the vernier calipers. Add one rod diameter to get $r$, the distance between the centres
        ($r = \qty{5.0}{mm}$).
  \item Join the rods at one end with the flexible link. The current goes along the top rod and
        back along the lower rod, so the two currents are opposite. Tare the balance.
  \item Set $I = \qty{4.0}{A}$, read $\Delta m$ quickly, and switch off. Repeat for 6.0, 8.0, 10.0
        and \qty{12.0}{A}, three times each. The same current goes through both rods, so
        $I_1 = I_2 = I$ and $I_1 I_2 = I^2$. Find $F = \Delta m\, g$ and plot $F$ against $I^2$.
  \item To test $F \propto 1/r$, keep $I$ fixed and change $r$ with spacers (small insulating
        blocks of known thickness).
  \item To see attraction, remove the flexible link. Join the two left ends to one supply
        terminal and the two right ends to the other. Now both rods carry current the same way
        (each carries about half of the supply current). The reading falls below zero: the lower
        rod is pulled up, towards the upper rod.
\end{enumerate}
Simple school demo: hang two long strips of aluminium foil side by side, about \qty{1}{cm} apart.
Pass a current along both strips the same way: they pull together. Connect them so the current
goes down one and up the other: they push apart.}
\results{\begin{sidebyside}[0.43]
    \centering\small
    \begin{tabular}{@{}cccc@{}}
      \toprule
      $I$ (\si{A}) & $I^2$ (\si{A^2}) & $\Delta m$ (\si{g}) & $F$ ($10^{-3}$ \si{N})\\
      \midrule
      4.0 & 16 & 0.032 & 0.31\\
      6.0 & 36 & 0.071 & 0.70\\
      8.0 & 64 & 0.127 & 1.24\\
      10.0 & 100 & 0.197 & 1.93\\
      12.0 & 144 & 0.285 & 2.79\\
      \bottomrule
    \end{tabular}
  \nextside
    \centering
    \begin{tikzpicture}
      \begin{axis}[hsc, width=0.97\linewidth, height=0.44\linewidth,
          xmin=0, xmax=160, ymin=0, ymax=3.2, xtick={0,40,...,160}, ytick={0,1,...,3},
          xlabel={$I^2$ (\si{A^2})}, ylabel={$F$ ($10^{-3}$ \si{N})}]
        \addplot[hscseriesA, only marks, mark=*, mark size=2pt]
          coordinates {(16,0.3136) (36,0.6958) (64,1.2446) (100,1.9306) (144,2.793)};
        \addplot[hscseriesB, dashed, line width=1pt, domain=0:155, samples=2, mark=none] {0.019376*x};
      \end{axis}
    \end{tikzpicture}
    \captionof{figure}{Force against current squared.}\label{ch2:fig:fvsi2}
  \end{sidebyside}
  Sample class results ($l = \qty{0.50}{m}$, $r = \qty{5.0e-3}{m}$, opposite currents). The
  reading rose: the lower rod was pushed down.}
\conclusion{The points (\cref{ch2:fig:fvsi2}) lie close to a straight line through the
  origin (0, 0), so $F \propto I^2$. Here $I_1 = I_2 = I$, so this agrees with $F \propto I_1 I_2$. Opposite
  currents repelled. Theory: $F = \left(\dfrac{\mu_0 l}{2\pi r}\right) I^2$, so the gradient is
  $\dfrac{\mu_0 l}{2\pi r}$. Using (0, 0) and (\qty{144}{A^2}, \qty{2.79e-3}{N}) on the line, the
  gradient is $\qty{2.79e-3}{N} \div \qty{144}{A^2} = \qty{1.94e-5}{N.A^{-2}}$. Rearrange to find
  $\mu_0$:
  \begin{align*}
    \text{gradient} &= \frac{\mu_0 l}{2\pi r} && \text{from the theory}\\
    2\pi r \times \text{gradient} &= \mu_0 l && \text{multiply both sides by $2\pi r$}\\
    \mu_0 &= \frac{2\pi r \times \text{gradient}}{l} && \text{divide both sides by $l$}
  \end{align*}
  \[ \mu_0 = \frac{2\pi \times \qty{5.0e-3}{m} \times \qty{1.94e-5}{N.A^{-2}}}{\qty{0.50}{m}}
  = \qty{1.22e-6}{N.A^{-2}}. \]
  This is about \qty{3}{\percent} below $4\pi \times 10^{-7} = \qty{1.26e-6}{N.A^{-2}}$. The value
  is a little low because the rods are not infinitely long: their ends make the force about
  \qty{1}{\percent} smaller. The rest of the difference is within the uncertainty in $r$ (about
  \qty{4}{\percent}). So the results support the formula.}
\significance{Ampère discovered this force in 1820. It defined the ampere from 1948 to 2019.}
\validity{\begin{itemize}
  \item \textbf{Validity:} the formula needs long, straight, parallel wires: use rods much longer
        than the gap ($l/r = 100$). Keep supply leads away, and use a light, flexible link.
  \item \textbf{Reliability:} repeat each current three times and average the readings.
  \item \textbf{Accuracy:} $r$ is small, so its error matters most (\qty{0.2}{mm} in \qty{5.0}{mm}
        is \qty{4}{\percent}): use vernier calipers and measure at both ends. Use short bursts, as
        hot rods bend.
  \item \textbf{Safety:} large currents make the rods hot (burns, fire): low voltage, short bursts.
\end{itemize}}
\end{experiment}

\begin{examtip}[breakable]
``Describe an investigation of the force on a current-carrying wire'' (3--4 marks). One mark for
each:
\begin{enumerate}
  \item a labelled diagram (magnet on a balance, wire clamped at \ang{90} to $B$, ammeter,
        rheostat);
  \item change $I$, and keep $l$, $B$ and $\theta$ the same;
  \item find $F = \Delta m\, g$ and plot $F$ against $I$: a straight line through the origin with
        gradient $lB$;
  \item turn the wire parallel to $B$: there is no force. Reverse $I$ or $B$: the force reverses.
\end{enumerate}
``Describe an investigation of the force between two parallel wires'' (about 4 marks). One mark
for each:
\begin{enumerate}
  \item a labelled diagram (supply, ammeter, rheostat, both wires or two hanging foil strips, the
        gap $r$);
  \item what you change, what you measure and what you keep the same;
  \item how you get $F$ ($F = \Delta m\, g$) and what you plot ($F$ against $I_1 I_2$);
  \item what you see (same direction: attract; opposite directions: repel).
\end{enumerate}
A ``hazard'' question wants the hazard \emph{and} a precaution.
\end{examtip}

\begin{quickcheck}
  \item A current balance reading rises by \qty{0.35}{g}. Find the force.
  \item In the parallel-rods experiment, $l = \qty{0.50}{m}$, $r = \qty{5.0}{mm}$ and the current
        is \qty{10.0}{A}. The rods are in series. Find the change in the balance reading,
        $\Delta m$, in grams.\label{ch2:qc:rods}
  \item Explain why $F$ against $I^2$ is a straight line. State its gradient.
  \item The rods are now connected so that both carry \qty{10.0}{A} in the same direction. Explain
        the new reading.
\end{quickcheck}

%% ==================================================================
\begin{summarybox}{\Cref{ch2:chap} in five lines}
\begin{itemize}
  \item A current makes a magnetic field: circles round a wire, $B = \dfrac{\mu_0 I}{2\pi r}$.
  \item A wire in an external field feels $F = lIB\sin\theta$: largest at \ang{90}, zero when
        parallel.
  \item Direction: right-hand palm rule (fingers $B$, thumb conventional $I$, palm $F$).
  \item Parallel wires: same direction attract, opposite directions repel;
        $\dfrac{F}{l} = \dfrac{\mu_0}{2\pi}\dfrac{I_1 I_2}{r}$.
  \item The forces are equal and opposite (Newton's third law), so one force value could define
        the ampere.
\end{itemize}
\end{summarybox}
